\documentclass[12pt, letterpaper]{article}

\usepackage[utf8]{inputenc}
\usepackage[T1]{fontenc}
\usepackage[spanish]{babel}
\usepackage{geometry}
\usepackage{amsmath, amssymb, amsthm}
\usepackage{booktabs}
\usepackage{array}
\usepackage{microtype}
\usepackage{setspace}
\usepackage{parskip}
\usepackage{titlesec}
\usepackage{fancyhdr}
\usepackage{enumitem}
\usepackage{bm}
\usepackage{mathtools}
\usepackage{tcolorbox}
\usepackage{hyperref}

\tcbuselibrary{skins, breakable}

\geometry{top=2.8cm, bottom=3.0cm, left=3.2cm, right=3.2cm}
\setstretch{1.20}

\pagestyle{fancy}
\fancyhf{}
\fancyhead[C]{\small\textsc{Econometria · Gujarati --- Ejercicio 11.1}}
\fancyfoot[C]{\small\thepage}
\renewcommand{\headrulewidth}{0.4pt}
\renewcommand{\footrulewidth}{0pt}

\titleformat{\section}
  {\large\bfseries\scshape}{\thesection.}{0.6em}{}
  [\vspace{-4pt}\rule{\linewidth}{0.4pt}]
\titleformat{\subsection}{\normalsize\bfseries}{\thesubsection.}{0.5em}{}
\titleformat{\subsubsection}{\normalsize\itshape}{}{0em}{}

\newtheorem{prop}{Proposicion}
\newtheorem{lema}{Lema}
\theoremstyle{remark}
\newtheorem{obs}{Observacion}

\newtcolorbox{verdict}[2]{
  colback=white, colframe=black,
  boxrule=0.5pt, arc=3pt,
  left=8pt, right=8pt, top=5pt, bottom=5pt,
  title={\textbf{Inciso #1 --- #2}},
  fonttitle=\small\bfseries,
  breakable
}

\newtcolorbox{clave}{
  colback=white, colframe=black,
  boxrule=0.5pt, arc=3pt,
  left=8pt, right=8pt, top=6pt, bottom=6pt,
  breakable
}

\newcommand{\E}{\mathbb{E}}
\newcommand{\Var}{\operatorname{Var}}
\newcommand{\Cov}{\operatorname{Cov}}
\newcommand{\plim}{\operatorname{plim}}

\begin{document}

\begin{center}
  {\LARGE\bfseries Verdadero, Falso o Incierto:}\\[6pt]
  {\LARGE\bfseries Siete Afirmaciones sobre Heteroscedasticidad}\\[12pt]
  {\large\itshape Justificacion formal con algebra matricial, teoria asintótica}\\
  {\large\itshape y analisis de casos}\\[14pt]
  \rule{0.55\linewidth}{0.6pt}\\[8pt]
  {\normalsize Ejercicio~11.1 --- \textit{Econometria}, Gujarati \& Porter, 5.a ed. por Emanuel Quintana Silva}\\[4pt]
  {\small\textsc{Con demostraciones del teorema de Gauss-Markov bajo varianza no constante,
  sesgo de los estimadores de varianza y prueba de White}}
\end{center}

\vspace{10pt}

\noindent\textbf{Nota preliminar.}\quad
El modelo de referencia es:
\begin{equation}
  Y_i = \beta_1 + \beta_2 X_{2i} + \cdots + \beta_k X_{ki} + u_i,
  \qquad i = 1, \ldots, n
  \label{eq:modelo}
\end{equation}
con el supuesto de heteroscedasticidad:
\begin{equation}
  \E[u_i] = 0, \quad \E[u_i^2] = \sigma_i^2, \quad \E[u_i u_j] = 0 \; (i \neq j)
  \label{eq:heterosc}
\end{equation}
donde $\sigma_i^2$ no es constante entre observaciones. En forma matricial,
la matriz de varianzas-covarianzas del vector de errores es:
\begin{equation}
  \E[\mathbf{u}\mathbf{u}'] = \boldsymbol\Omega =
  \begin{pmatrix}
    \sigma_1^2 & 0 & \cdots & 0\\
    0 & \sigma_2^2 & \cdots & 0\\
    \vdots & \vdots & \ddots & \vdots\\
    0 & 0 & \cdots & \sigma_n^2
  \end{pmatrix}
  \neq \sigma^2 \mathbf{I}_n
  \label{eq:Omega}
\end{equation}

\bigskip

%% ── Inciso (a) ────────────────────────────────────────────────────────────
\begin{verdict}{(a)}{FALSO}
  En presencia de heteroscedasticidad, los estimadores de MCO son sesgados e
  ineficientes.
\end{verdict}

\subsection*{Justificacion}

El estimador MCO es $\hat{\boldsymbol\beta} = (\mathbf{X'X})^{-1}\mathbf{X'y}$.
Sustituyendo $\mathbf{y} = \mathbf{X}\boldsymbol\beta + \mathbf{u}$:
\begin{equation}
  \hat{\boldsymbol\beta}
    = \boldsymbol\beta + (\mathbf{X'X})^{-1}\mathbf{X'u}
  \label{eq:beta_descomp}
\end{equation}
Tomando esperanza bajo exogeneidad ($\E[\mathbf{X'u}] = \mathbf{0}$):
\begin{equation}
  \E[\hat{\boldsymbol\beta}] = \boldsymbol\beta
  \label{eq:insesgado}
\end{equation}

El estimador MCO es \textbf{insesgado} independientemente de si hay
heteroscedasticidad. La insesgadez no requiere homoscedasticidad: solo
requiere exogeneidad de los regresores ($\E[u_i | \mathbf{X}] = 0$),
que no se ve afectada por la varianza no constante.

Lo que se pierde bajo heteroscedasticidad es la \textbf{eficiencia}: el
estimador MCO ya no es MELI (Mejor Estimador Lineal Insesgado). La varianza
del estimador MCO bajo heteroscedasticidad es:
\begin{equation}
  \Var(\hat{\boldsymbol\beta}) = (\mathbf{X'X})^{-1}\mathbf{X'}\boldsymbol\Omega\mathbf{X}(\mathbf{X'X})^{-1}
  \label{eq:var_heterosc}
\end{equation}
que es mayor que la varianza del estimador de Minimos Cuadrados Generalizados
(MCG):
\begin{equation}
  \Var(\hat{\boldsymbol\beta}^{\text{MCG}}) = (\mathbf{X'}\boldsymbol\Omega^{-1}\mathbf{X})^{-1}
  \leq \Var(\hat{\boldsymbol\beta}^{\text{MCO}})
  \label{eq:MCG_eficiente}
\end{equation}
por el teorema de Gauss-Markov generalizado.

\textbf{Conclusion}: \textsc{Falso}. MCO sigue siendo insesgado y consistente.
Solo pierde eficiencia: ya no es MELI.

\bigskip

%% ── Inciso (b) ────────────────────────────────────────────────────────────
\begin{verdict}{(b)}{VERDADERO}
  Si hay heteroscedasticidad, las pruebas convencionales $t$ y $F$ son invalidas.
\end{verdict}

\subsection*{Justificacion}

Las pruebas convencionales $t$ y $F$ se construyen usando el estimador de
varianza homoscedastico:
\begin{equation}
  \hat\sigma^2_{\text{conv}} = \frac{\hat{\mathbf{u}}'\hat{\mathbf{u}}}{n - k}
  \label{eq:sigma2_conv}
\end{equation}
y la formula de varianza que asume $\boldsymbol\Omega = \sigma^2\mathbf{I}$:
\begin{equation}
  \widehat{\Var}(\hat\beta_j)_{\text{conv}} = \hat\sigma^2[(\mathbf{X'X})^{-1}]_{jj}
  \label{eq:var_conv}
\end{equation}

Bajo heteroscedasticidad, la verdadera varianza de $\hat\beta_j$ es el
elemento $j,j$ de~(\ref{eq:var_heterosc}), que en general difiere de~(\ref{eq:var_conv}).
Por tanto, el estimador convencional~(\ref{eq:var_conv}) es un estimador
\textbf{sesgado} de la verdadera varianza, lo que implica que:

\begin{itemize}[leftmargin=1.6em, itemsep=3pt]
  \item El estadistico $t_j = \hat\beta_j / \text{ee}_{\text{conv}}(\hat\beta_j)$
    no sigue la distribucion $t_{n-k}$ bajo $H_0$, porque el denominador es incorrecto.
  \item El estadistico $F$ tampoco sigue la distribucion $F_{k-1, n-k}$.
\end{itemize}

El remedio es usar errores estandar robustos a la heteroscedasticidad
(estimadores de White, 1980):
\begin{equation}
  \widehat{\Var}^{\text{White}}(\hat{\boldsymbol\beta})
    = (\mathbf{X'X})^{-1}
      \left(\sum_i \hat u_i^2 \mathbf{x}_i\mathbf{x}_i'\right)
      (\mathbf{X'X})^{-1}
  \label{eq:White_var}
\end{equation}
que es un estimador consistente de~(\ref{eq:var_heterosc}).

\textbf{Conclusion}: \textsc{Verdadero}. Las pruebas $t$ y $F$ convencionales
son invalidas bajo heteroscedasticidad porque los errores estandar usados en
su construccion son sesgados.

\bigskip

%% ── Inciso (c) ────────────────────────────────────────────────────────────
\begin{verdict}{(c)}{FALSO}
  En presencia de heteroscedasticidad, el metodo de MCO habitual siempre
  sobreestima los errores estandar de los estimadores.
\end{verdict}

\subsection*{Justificacion}

El sesgo del estimador convencional de varianza depende de la forma de la
heteroscedasticidad y de la estructura de los datos. Para ilustrarlo en
el caso simple de regresion con un regresor:

\begin{equation}
  \Var(\hat\beta_2)_{\text{verdadera}}
    = \frac{\sum_i \sigma_i^2 x_i^2}{(\sum_i x_i^2)^2}
  \label{eq:var_verdadera}
\end{equation}
\begin{equation}
  \E[\widehat{\Var}(\hat\beta_2)_{\text{conv}}]
    = \sigma^2 [(\mathbf{X'X})^{-1}]_{22}
  \label{eq:var_conv_esp}
\end{equation}

donde $\sigma^2$ es algun promedio de las $\sigma_i^2$. El signo del sesgo
$\E[\widehat{\Var}_{\text{conv}}] - \Var_{\text{verdadera}}$ depende de si
las varianzas $\sigma_i^2$ son mayores o menores para las observaciones con
$x_i^2$ grande:

\begin{itemize}[leftmargin=1.6em, itemsep=4pt]
  \item Si $\sigma_i^2$ es grande cuando $x_i^2$ es grande (heteroscedasticidad
    ``tipica'', como en datos de ingresos), la varianza verdadera puede ser
    mayor que la convencional: MCO \emph{subestima} el error estandar.
  \item Si $\sigma_i^2$ es grande cuando $x_i^2$ es pequeno, la varianza
    convencional puede exceder la verdadera: MCO \emph{sobreestima} el error.
\end{itemize}

La afirmacion de que MCO ``siempre sobreestima'' es incorrecta: el sesgo puede
ir en cualquier direccion.

\textbf{Conclusion}: \textsc{Falso}. El sesgo del estimador convencional de
varianza puede ser positivo o negativo; no hay una direccion unica.

\bigskip

%% ── Inciso (d) ────────────────────────────────────────────────────────────
\begin{verdict}{(d)}{INCIERTO}
  Si los residuales estimados mediante una regresion por MCO exhiben un patron
  sistematico, significa que hay heteroscedasticidad en los datos.
\end{verdict}

\subsection*{Justificacion}

Un patron sistematico en los residuos $\hat u_i$ puede originarse en multiples
fuentes, no solo en heteroscedasticidad:

\begin{enumerate}[leftmargin=1.8em, itemsep=5pt]
  \item \textbf{Heteroscedasticidad}: si $\sigma_i^2 = f(X_i)$, los residuos
    mostraran mayor dispersion para valores grandes de $X_i$ (patron de abanico).

  \item \textbf{Autocorrelacion}: si los errores estan correlacionados en el
    tiempo ($\Cov(u_t, u_{t-s}) \neq 0$), los residuos mostraran runs o ciclos
    sistematicos que se pueden confundir con heteroscedasticidad.

  \item \textbf{Forma funcional incorrecta}: si el verdadero modelo es
    $Y = \beta_1 + \beta_2 X + \beta_3 X^2 + u$ pero se estima solo la parte
    lineal, los residuos tendran una curvatura sistematica.

  \item \textbf{Variable omitida}: si se omite una variable relevante $Z$,
    los residuos capturan $\beta_Z Z_i + u_i$, lo que crea patron si $Z$ varia
    sistematicamente con $X$.
\end{enumerate}

La inspeccion visual de residuos es un primer diagnostico util, pero no
es concluyente. Para distinguir entre estas causas se necesitan pruebas
formales (Breusch-Pagan, White, Durbin-Watson, RESET, etc.).

\textbf{Conclusion}: \textsc{Incierto}. Un patron en los residuos es necesario
pero no suficiente para concluir heteroscedasticidad; puede deberse a otras
formas de mala especificacion.

\bigskip

%% ── Inciso (e) ────────────────────────────────────────────────────────────
\begin{verdict}{(e)}{FALSO}
  No hay una prueba general de heteroscedasticidad que no este basada en algun
  supuesto acerca de cual variable esta correlacionada con el termino de error.
\end{verdict}

\subsection*{Justificacion}

La \textbf{prueba de White} (1980) es una prueba general que no requiere
supuestos previos sobre cual variable causa la heteroscedasticidad.
El procedimiento es:

\begin{enumerate}[leftmargin=1.8em, itemsep=4pt]
  \item Estimar el modelo~(\ref{eq:modelo}) por MCO y obtener los residuos $\hat u_i$.
  \item Regresar $\hat u_i^2$ sobre todos los regresores, sus cuadrados y sus
    productos cruzados:
    \begin{equation}
      \hat u_i^2 = \alpha_0 + \sum_j \alpha_j X_{ji}
                 + \sum_j \alpha_{jj} X_{ji}^2
                 + \sum_{j<k} \alpha_{jk} X_{ji} X_{ki}
                 + v_i
      \label{eq:white_reg}
    \end{equation}
  \item El estadistico de prueba es $nR^2_{\hat u^2}$, donde $R^2_{\hat u^2}$
    es el $R^2$ de la regresion auxiliar~(\ref{eq:white_reg}):
    \begin{equation}
      W = nR^2_{\hat u^2} \xrightarrow{d} \chi^2_p
      \label{eq:white_stat}
    \end{equation}
    donde $p$ es el numero de regresores en~(\ref{eq:white_reg}) (sin el intercepto).
\end{enumerate}

La prueba de White no supone ninguna forma funcional especifica para
$\sigma_i^2 = f(\cdot)$: detecta cualquier forma de heteroscedasticidad que
este relacionada con los regresores, sus cuadrados o sus productos cruzados.

\textbf{Conclusion}: \textsc{Falso}. La prueba de White es una prueba general
que no requiere supuesto previo sobre cual variable genera la varianza no constante.

\bigskip

%% ── Inciso (f) ────────────────────────────────────────────────────────────
\begin{verdict}{(f)}{VERDADERO (con matices)}
  Si el modelo de regresion esta mal especificado (por ejemplo, si se omitio
  una variable importante), los residuos de MCO mostraran un patron claramente
  distinguible.
\end{verdict}

\subsection*{Justificacion}

Si el verdadero modelo es:
\begin{equation}
  Y_i = \beta_1 + \beta_2 X_{2i} + \beta_3 X_{3i} + u_i
  \label{eq:verdadero}
\end{equation}
pero se estima omitiendo $X_3$:
\begin{equation}
  Y_i = \alpha_1 + \alpha_2 X_{2i} + e_i, \qquad e_i = u_i + \beta_3 X_{3i}
  \label{eq:reducido}
\end{equation}
El termino de error compuesto $e_i = u_i + \beta_3 X_{3i}$ contiene la
influencia sistematica de $X_{3i}$. Si $X_3$ varia sistematicamente con
$X_2$, los residuos del modelo reducido mostraran un patron vinculado a
$X_2$: al graficar $\hat e_i$ contra $X_{2i}$, aparecera una tendencia
que no deberia existir si el modelo estuviera bien especificado.

Analogamente, si la forma funcional es incorrecta (modelo lineal cuando
la verdadera relacion es cuadratica), los residuos tendran curvatura
sistematica al graficarlos contra $X$.

El matiz es que el patron puede no ser siempre ``claramente distinguible''
en muestras pequenas con mucho ruido, pero en principio la mala especificacion
introduce componentes no aleatorios en los residuos.

\textbf{Conclusion}: \textsc{Verdadero}. La mala especificacion introduce
componentes sistematicos en los residuos que producen patrones detectables.

\bigskip

%% ── Inciso (g) ────────────────────────────────────────────────────────────
\begin{verdict}{(g)}{VERDADERO}
  Si una regresora con varianza no constante se omite incorrectamente de un
  modelo, los residuos (MCO) seran heteroscedasticos.
\end{verdict}

\subsection*{Justificacion}

Suponga que el verdadero modelo es:
\begin{equation}
  Y_i = \beta_1 + \beta_2 X_{2i} + \beta_3 X_{3i} + u_i,
  \quad \E[u_i^2] = \sigma^2 \text{ (homosc.)}
  \label{eq:verdadero_g}
\end{equation}
y que $X_{3i}$ tiene varianza condicional no constante en el sentido de que
$\Var(X_{3i} | X_{2i}) = h(X_{2i})$ no es constante. Si se omite $X_3$ y
se estima:
\begin{equation}
  Y_i = \alpha_1 + \alpha_2 X_{2i} + v_i, \quad v_i = u_i + \beta_3 X_{3i}
  \label{eq:reducido_g}
\end{equation}
la varianza del error compuesto es:
\begin{equation}
  \Var(v_i | X_{2i})
    = \Var(u_i | X_{2i}) + \beta_3^2 \Var(X_{3i} | X_{2i})
    = \sigma^2 + \beta_3^2 h(X_{2i})
  \label{eq:var_vi}
\end{equation}
Como $h(X_{2i})$ no es constante, $\Var(v_i | X_{2i})$ tampoco lo es:
el modelo reducido tiene heteroscedasticidad aunque el verdadero modelo sea
homoscedastico. Los residuos del modelo reducido exhibiran varianza no
constante inducida por la omision.

Este resultado subraya una fuente frecuentemente pasada por alto de
heteroscedasticidad aparente: la mala especificacion del modelo puede
generar residuos heteroscedasticos que no reflejan varianza no constante
del proceso verdadero, sino la omision de una variable relevante con
variabilidad no uniforme.

\textbf{Conclusion}: \textsc{Verdadero}. Al omitir una variable con varianza
condicional no constante, el error del modelo reducido hereda esa
heteroscedasticidad.

\bigskip

%% ── Tabla resumen ─────────────────────────────────────────────────────────
\section{Tabla resumen}

\begin{table}[ht]
\centering
\caption{Veredictos del ejercicio~11.1}
\label{tab:resumen}
\small
\begin{tabular}{clp{9cm}}
\toprule
\textbf{Inciso} & \textbf{Veredicto} & \textbf{Razon principal} \\
\midrule
(a) & Falso
  & MCO sigue siendo insesgado y consistente; solo pierde eficiencia (MELI). \\[4pt]
(b) & Verdadero
  & Las varianzas convencionales son estimadores sesgados bajo heterosc.; los estadisticos $t$ y $F$ no siguen sus distribuciones nominales. \\[4pt]
(c) & Falso
  & El sesgo del estimador convencional de varianza puede ser positivo o negativo segun la forma de la heterosc. \\[4pt]
(d) & Incierto
  & Patron en residuos puede deberse tambien a autocorrelacion, forma funcional incorrecta o variable omitida. \\[4pt]
(e) & Falso
  & La prueba de White (1980) es general: no supone que variable causa la heterosc. \\[4pt]
(f) & Verdadero
  & Mala especificacion introduce componentes sistematicos en los residuos. \\[4pt]
(g) & Verdadero
  & La varianza del error compuesto hereda la heterogeneidad de la variable omitida. \\
\bottomrule
\end{tabular}
\end{table}

\appendix

\section{Demostracion del sesgo del estimador convencional de varianza}

Bajo heteroscedasticidad, el valor esperado del estimador convencional
$\hat\sigma^2 = \hat{\mathbf{u}}'\hat{\mathbf{u}}/(n-k)$ es:

\begin{equation}
  \E[\hat\sigma^2]
    = \frac{1}{n-k}\E[\hat{\mathbf{u}}'\hat{\mathbf{u}}]
    = \frac{1}{n-k}\operatorname{tr}(\mathbf{M}\boldsymbol\Omega)
  \label{eq:E_sigma2}
\end{equation}

donde $\mathbf{M} = \mathbf{I} - \mathbf{X}(\mathbf{X'X})^{-1}\mathbf{X}'$
es la matriz de proyeccion ortogonal sobre el complemento del espacio columna
de $\mathbf{X}$. Bajo homoscedasticidad ($\boldsymbol\Omega = \sigma^2\mathbf{I}$):

\begin{equation}
  \E[\hat\sigma^2] = \frac{\sigma^2}{n-k}\operatorname{tr}(\mathbf{M}) = \sigma^2
  \label{eq:E_sigma2_homosc}
\end{equation}

Bajo heteroscedasticidad, $\operatorname{tr}(\mathbf{M}\boldsymbol\Omega) \neq (n-k)\sigma^2_{\text{promedio}}$
en general, y el estimador puede ser mayor o menor que la varianza promedio.

\section{Por que el estimador de White es consistente}

El estimador de White~(\ref{eq:White_var}) es consistente porque:

\begin{equation}
  \frac{1}{n}\sum_i \hat u_i^2 \mathbf{x}_i\mathbf{x}_i'
  \xrightarrow{p} \frac{1}{n}\sum_i \sigma_i^2 \mathbf{x}_i\mathbf{x}_i'
  \label{eq:white_consistent}
\end{equation}

bajo condiciones de regularidad (Ley de los Grandes Numeros). La convergencia
usa el hecho de que $\hat u_i = u_i - (\hat{\boldsymbol\beta} - \boldsymbol\beta)'\mathbf{x}_i$
y que $\hat{\boldsymbol\beta} \xrightarrow{p} \boldsymbol\beta$, lo que hace
que $\hat u_i \xrightarrow{p} u_i$ asintoticamente.

\vspace{14pt}
\noindent\rule{\linewidth}{0.4pt}

\noindent{\small\textit{Nota.} La distincion entre insesgadez y eficiencia del
inciso~(a) es fundamental: el teorema de Gauss-Markov requiere homoscedasticidad
para garantizar eficiencia (minima varianza), pero la insesgadez solo requiere
exogeneidad ($\E[u_i|\mathbf{X}] = 0$). Estos son supuestos distintos y la
violacion de uno no implica la violacion del otro.}

\end{document}
